\chapter{Árboles de juego}

\section{Modelo formal de un juego finito}

\begin{definition}[Juego bipersonal finito, determinista y de suma cero]
Un juego bipersonal finito, determinista y de suma cero se modela mediante la tupla
\[
G=(\States,s_0,\MAX,\MIN,\Player,\Legal,\Succ,\Term,\Utility).
\]
Sus componentes son:
\begin{itemize}
\item $\States$: conjunto de estados posibles;
\item $s_0\in\States$: estado inicial;
\item $\MAX$ y $\MIN$: los dos jugadores;
\item $\Player(s)$: jugador que mueve en el estado $s$;
\item $\Legal(s)$: conjunto de movimientos legales en $s$;
\item $\Succ(s,a)$: estado sucesor al aplicar el movimiento $a$ en $s$;
\item $\Term(s)$: predicado que indica si $s$ es terminal;
\item $\Utility(s)$: valor numérico de un estado terminal desde el punto de vista de $\MAX$.
\end{itemize}
\end{definition}

\begin{intuitionbox}
``Determinista'' significa que una jugada produce siempre el mismo resultado. ``Suma cero'' significa que lo que favorece a un jugador perjudica al otro. Si MAX gana $1$, MIN pierde $1$; si MAX pierde $1$, MIN gana $1$.
\end{intuitionbox}

\section{Construcción del árbol}

\begin{definition}[Árbol de juego]
El árbol de juego asociado a un estado $s$ es el árbol enraizado en $s$ tal que:
\begin{enumerate}
\item la raíz es $s$;
\item si $u$ es un nodo no terminal, sus hijos son exactamente los estados de $\Children(u)$;
\item si $u$ es terminal, no tiene hijos.
\end{enumerate}
\end{definition}

\begin{formalbox}
La regla de construcción puede escribirse así:
\[
\Children(s)=\{\Succ(s,a):a\in\Legal(s)\}.
\]
Si $\Term(s)=\true$, entonces $\Legal(s)=\varnothing$ y $\Children(s)=\varnothing$.
\end{formalbox}

\section{Niveles, profundidad y factor de ramificación}

\begin{definition}[Profundidad]
La profundidad de un nodo es la cantidad de aristas que hay desde la raíz hasta ese nodo. La raíz tiene profundidad $0$.
\end{definition}

\begin{definition}[Factor de ramificación]
El factor de ramificación de un estado $s$ es el número de hijos de $s$:
\[
 b(s)=|\Children(s)|.
\]
Cuando todos los nodos tienen aproximadamente $b$ hijos, se dice informalmente que el árbol tiene factor de ramificación $b$.
\end{definition}

\section{Por qué los árboles crecen tan rápido}

Si cada estado tiene $b$ hijos y exploramos hasta profundidad $d$, el número de hojas es aproximadamente
\[
 b^d.
\]

\begin{example}[Crecimiento con $b=3$]
Si $b=3$ y $d=4$, hay $3^4=81$ hojas. Esto coincide con el tipo de árbol usado en los ejemplos clásicos de análisis de alpha--beta: pocas ramas por nodo, pero explosión exponencial al crecer la profundidad.
\end{example}

\begin{summarybox}
El árbol de juego convierte un problema estratégico en un problema de búsqueda. Minimax asignará valores a los nodos; Alpha--Beta intentará evitar visitar nodos cuyo valor no puede cambiar la decisión final.
\end{summarybox}
